\documentclass[aps,prd,preprint,notitlepage,nofootinbib,longbibliography]{revtex4-2}

\usepackage{amsmath,amssymb,bm}
\usepackage{graphicx}
\usepackage{booktabs}
\usepackage{microtype}
\usepackage{hyperref}
\usepackage{enumitem}
\usepackage{siunitx}

\hypersetup{
  colorlinks=true,
  linkcolor=blue,
  citecolor=blue,
  urlcolor=blue,
  pdftitle={Elastic Worldtube Memory as a Sandbox Model for Interbraid Backtug},
  pdfauthor={Nathan McKnight}
}

\newcommand{\HsH}{H(s)H}
\newcommand{\kb}{k_{\rm b}}
\newcommand{\lel}{\ell_{\rm el}}
\newcommand{\lI}{\ell_{\rm I}}

\begin{document}

\title{Elastic Worldtube Memory as a Sandbox Model for Interbraid Backtug:\\
A Constitutive Transition and an Independent Dispersion Test}

\author{Nathan McKnight}
\affiliation{Independent Researcher, Columbus, Ohio, USA}
\thanks{This manuscript develops a sandbox construction within the SAT$\rightarrow$\HsH{} research program. Algebraic and numerical checks were assisted by an AI research instance designated Mercer.}

\date{October 2, 2026}

\begin{abstract}
SAT and its Hyperhelical Worldtube successor, \HsH{}, explore the consequences of treating particle histories as finite, physically structured four-dimensional worldtubes rather than as structureless bookkeeping curves. A recurring internal question is whether a localized interaction can leave a mechanically persistent relation along an extended worldtube, producing what the project calls \emph{Interbraid backtug}. This paper does not assume such a force exists. Instead, it constructs a minimal continuum null model and asks what would follow if an extended carrier possessed ordinary tension, bending stiffness, and a local pair-restoring response. For a relative deformation $u(s)$ the quadratic functional
\begin{equation*}
E[u]=\frac12\int ds\left[B(u'')^2+T(u')^2+\kb u^2\right]
\end{equation*}
leads to the static operator $B\partial_s^4-T\partial_s^2+\kb$. Its Green function exhibits a sharp response-class transition controlled by the dimensionless number $\chi=4B\kb/T^2$. For $\chi<1$, a localized interaction produces a sum of monotonic exponentials. At $\chi=1$ the characteristic roots merge. For $\chi>1$, the stored deformation is an exponentially damped spatial oscillation, so the induced restoring contribution alternates sign along the carrier. The same parameters independently predict the undamped dispersion relation $\omega^2=(Bq^4+Tq^2+\kb)/\mu$. Consequently the static profile and dynamic spectrum provide a two-channel falsification test of the model. The construction does not establish gravity, electromagnetism, quantum entanglement, retrocausality, or a physical worldtube medium. It supplies a calculable sandbox mechanism by which interaction history could be stored as present geometry, and a precise criterion by which that mechanism can fail.
\end{abstract}

\maketitle

\begin{center}
\fbox{\parbox{0.93\linewidth}{\textbf{SANDBOX STATUS.} This manuscript is a model-development exercise. SAT/\HsH{} terminology such as ``Interbraid'' and ``Electrogravity'' is project-internal nomenclature, not an assertion of established forces. The only claims promoted here are mathematical consequences of the explicitly stated continuum model.}}
\end{center}

\section{Introduction}

A persistent theme in the SAT$\rightarrow$\HsH{} archive is deliberately simple: begin with the standard four-dimensional representation of kinematics, then ask what changes if the worldline is replaced by a finite, physically structured worldtube. Two archival questions motivate the present calculation. First, an older SAT note asks whether long interaction history within a bound system could leave a present dynamical difference between long-associated and newly arriving bodies \cite{gravitytwist}. Second, a later \HsH{} conversation poses the stripped-down geometric problem directly: if two extended particle histories interact at one location along their four-dimensional extent, could the physical carrier preserve some consequence of that encounter away from the interaction point \cite{forcesacrosstime}?

Neither question is answered merely by drawing a block universe. Standard relativistic spacetime geometry does not, by itself, add a new cross-temporal force. A physical mechanism must be specified. The purpose of this paper is therefore narrower: construct the least exotic extended-carrier model that can store a localized interaction as present deformation, derive its consequences completely enough to be falsifiable, and identify which consequences are geometry, which are constitutive assumptions, and which remain conjectural.

The model is intentionally a null model. It borrows the mathematics of a tension-bearing, bendable continuum with a local restoring term. If a future \HsH{} worldtube cannot support anything analogous to these degrees of freedom, the model should be rejected. If it can, the model supplies an explicit answer to a useful sub-question: what forms of persistent ``backtug'' arise from ordinary continuum mechanics before any new force law is invented?

\section{Provenance and modeling commitments}

Table~\ref{tab:provenance} separates the source material from the new construction. This is essential because the archive contains exploratory physical interpretations as well as later corrections and abandoned branches.

\begin{table}[h]
\caption{Boundary between source material and the present sandbox derivation.}
\label{tab:provenance}
\begin{tabular}{@{}ll@{}}
\toprule
Category & Content \\
\midrule
Source-derived question & \parbox[t]{0.69\linewidth}{Treat particle histories as physically extended worldtubes and ask whether a localized encounter can leave a persistent relation along their extent \cite{gravitytwist,forcesacrosstime}.} \\[0.7em]
Source-derived program & \parbox[t]{0.69\linewidth}{Explore history-dependent, direct worldtube--worldtube mechanics separately from the shared surface/medium response called Electrogravity in current project terminology.} \\[0.7em]
New assumption here & \parbox[t]{0.69\linewidth}{Relative worldtube deformation admits a quadratic continuum energy with bending modulus $B$, tension $T$, and local restoring stiffness $\kb$.} \\[0.7em]
New derivation here & \parbox[t]{0.69\linewidth}{Exact static response classes, the critical ratio $\chi=4B\kb/T^2$, persistence lengths, damped spatial oscillations, and the independent dynamic dispersion relation.} \\[0.7em]
Not established here & \parbox[t]{0.69\linewidth}{Gravity, electromagnetism, the strong interaction, quantum entanglement, a retrocausal force, a particular particle identity, or an empirical anomaly.} \\
\bottomrule
\end{tabular}
\end{table}

The carrier coordinate $s$ below is an intrinsic coordinate along the reduced extended object. The static calculation is a configuration-space model. It does not by itself specify a causal update law in laboratory time, and therefore does not by itself imply retrocausality. A dynamic completion is introduced later only as an additional constitutive hypothesis.

\section{Minimal elastic model}

Let $u(s)$ denote a small relative displacement of two coupled carriers, or equivalently a reduced deformation coordinate for one carrier relative to an interaction channel. The minimal positive quadratic energy retaining curvature, slope, and local restoring costs is
\begin{equation}
E[u]=\frac12\int_{-\infty}^{\infty} ds\left[B\left(u''\right)^2+T\left(u'\right)^2+\kb u^2\right]-\int ds\, f(s)u(s).
\label{eq:energy}
\end{equation}
Here $B>0$ is a bending modulus, $T>0$ is an effective longitudinal tension, and $\kb>0$ is a local pair-restoring coefficient. For $[s]=[u]=L$,
\begin{equation}
[B]=F L^2,\qquad [T]=F,\qquad [\kb]=F/L^2,
\end{equation}
so every term in Eq.~\eqref{eq:energy} has dimensions of energy.

Variation gives
\begin{equation}
B u''''-T u''+\kb u=f(s).
\label{eq:static}
\end{equation}
A localized encounter is represented by the benchmark source
\begin{equation}
f(s)=F_0\delta(s).
\label{eq:delta}
\end{equation}
This source is not meant to model an actual particle collision in detail. It is a probe of the carrier response.

The Fourier-space Green function is
\begin{equation}
\widetilde G(q)=\frac{1}{Bq^4+Tq^2+\kb},
\label{eq:fouriergreen}
\end{equation}
and $u(s)=F_0G(s)$.

Two natural lengths appear immediately:
\begin{equation}
\lel=\sqrt{\frac{B}{T}},\qquad
\lI=\sqrt{\frac{T}{\kb}}.
\label{eq:lengths}
\end{equation}
The first is the familiar bending--tension crossover length. The second is the persistence length of a tension-bearing carrier against the local restoring channel. Their ratio generates the dimensionless control number
\begin{equation}
\boxed{\chi=\frac{4B\kb}{T^2}=4\left(\frac{\lel}{\lI}\right)^2.}
\label{eq:chi}
\end{equation}
Using $x=s/\lel$, the homogeneous static equation reduces to
\begin{equation}
\frac{d^4u}{dx^4}-\frac{d^2u}{dx^2}+\frac{\chi}{4}u=0,
\label{eq:dimensionless}
\end{equation}
showing that the response class is controlled by a single dimensionless number.

\section{Static response and the constitutive transition}

Seeking $u\propto e^{-rs}$ gives
\begin{equation}
Br^4-Tr^2+\kb=0.
\end{equation}
With $y=r^2$,
\begin{equation}
y_{\pm}=\frac{T\pm\sqrt{T^2-4B\kb}}{2B}.
\label{eq:yroots}
\end{equation}
The discriminant
\begin{equation}
\Delta=T^2-4B\kb=T^2(1-\chi)
\end{equation}
creates three response classes.

\subsection{Subcritical regime: $\chi<1$}

For $\chi<1$, both $y_{\pm}$ are positive. Define $r_\pm=\sqrt{y_\pm}$ with $r_+>r_->0$. The exact even Green function is
\begin{equation}
G(s)=\frac{1}{2B(r_+^2-r_-^2)}
\left[
\frac{e^{-r_-|s|}}{r_-}-\frac{e^{-r_+|s|}}{r_+}
\right].
\label{eq:greensubcritical}
\end{equation}
The response is monotonic and contains two decay lengths, $r_+^{-1}$ and $r_-^{-1}$. In the weak-bending limit,
\begin{equation}
r_+^{-1}\sim\sqrt{\frac{B}{T}}=\lel,
\qquad
r_-^{-1}\sim\sqrt{\frac{T}{\kb}}=\lI.
\end{equation}
Thus the same static response cleanly separates a short bending boundary layer from a longer interaction-persistence scale.

If bending is removed entirely, Eq.~\eqref{eq:static} becomes
\begin{equation}
-Tu''+\kb u=F_0\delta(s),
\end{equation}
with
\begin{equation}
G_{B=0}(s)=\frac{1}{2\sqrt{T\kb}}
\exp\left(-\frac{|s|}{\lI}\right).
\label{eq:tensiongreen}
\end{equation}
This limit makes the role of $\lI$ transparent. It also exposes a failure of tension alone: as $\kb\rightarrow0$, the localized static solution ceases to be well behaved on an infinite carrier. Persistent localized deformation therefore requires a restoring channel, boundary conditions, topological trapping, or some other mechanism beyond bare tension.

\subsection{Critical regime: $\chi=1$}

At
\begin{equation}
\boxed{T_c=2\sqrt{B\kb}},
\label{eq:Tcrit}
\end{equation}
the roots merge at
\begin{equation}
r_0^2=\frac{T}{2B}.
\end{equation}
The Green function becomes
\begin{equation}
\boxed{
G_{\rm crit}(s)=
\frac{(1+r_0|s|)e^{-r_0|s|}}{4Br_0^3}.
}
\label{eq:greencritical}
\end{equation}
In length-language the critical condition is
\begin{equation}
\boxed{\lI=2\lel.}
\label{eq:lengthcritical}
\end{equation}
This is a constitutive response boundary, not a thermodynamic phase transition.

\subsection{Supercritical regime: $\chi>1$}

For $\chi>1$, the squared characteristic roots are complex conjugates. Write the decaying roots as
\begin{equation}
r=\alpha\pm i\beta,
\end{equation}
where
\begin{align}
\alpha^2&=\frac12\left[\sqrt{\frac{\kb}{B}}+\frac{T}{2B}\right],\\
\beta^2&=\frac12\left[\sqrt{\frac{\kb}{B}}-\frac{T}{2B}\right].
\label{eq:alphabeta}
\end{align}
The exact even Green function may be written
\begin{equation}
\boxed{
G(s)=\frac{e^{-\alpha|s|}}{4B\alpha(\alpha^2+\beta^2)}
\left[
\cos(\beta|s|)+\frac{\alpha}{\beta}\sin(\beta|s|)
\right].
}
\label{eq:greensupercritical}
\end{equation}
The response is now a damped spatial oscillation. A local restoring contribution proportional to $-\kb u(s)$ therefore alternates sign along the carrier. In project language this supplies a concrete, non-mystical meaning for alternating ``backtug'' and ``forwardtug'': they are opposite phases of one underdamped spatial response, not two separately postulated forces.

The approximate separation between sign reversals is set by
\begin{equation}
\Delta s_{\rm flip}\sim\frac{\pi}{\beta}.
\label{eq:flip}
\end{equation}
Figure~\ref{fig:green} shows the three static response classes for a dimensionless benchmark with $B=T=1$ and $\kb=\chi/4$.

\begin{figure}[t]
\includegraphics[width=0.94\linewidth]{green_profiles.png}
\caption{Normalized static Green functions for subcritical, critical, and supercritical response. The supercritical curve crosses zero and develops a damped oscillatory tail. Parameters are dimensionless with $B=T=1$ and $\kb=\chi/4$. The figure is a numerical evaluation of Eq.~\eqref{eq:fouriergreen}, not a fit.}
\label{fig:green}
\end{figure}

\section{Independent dynamic cross-check}

A static stored deformation is more credible if the same constitutive parameters also govern a consistent dynamic response. Introduce a mass per carrier length $\mu$ and optional linear damping $\zeta$:
\begin{equation}
\mu u_{\tau\tau}+\zeta u_\tau+B u_{ssss}-T u_{ss}+\kb u=f(s,\tau).
\label{eq:dynamic}
\end{equation}
Here $\tau$ is an evolution parameter for the reduced carrier model. This equation is an additional assumption; the static model alone does not provide it.

For an undamped plane mode $u\propto e^{i(qs-\omega\tau)}$,
\begin{equation}
\boxed{
\omega^2(q)=\frac{Bq^4+Tq^2+\kb}{\mu}.
}
\label{eq:dispersion}
\end{equation}
The spectrum contains three regimes:
\begin{align}
q\rightarrow0 &: \quad \omega\rightarrow\sqrt{\kb/\mu},\\
q\lel\sim1 &: \quad \text{tension and bending compete},\\
q\lel\gg1 &: \quad \omega\sim\sqrt{B/\mu}\,q^2.
\end{align}
With
\begin{equation}
\widehat q=q\lel,
\qquad
\widehat\omega=\omega\frac{\sqrt{B\mu}}{T},
\end{equation}
Eq.~\eqref{eq:dispersion} becomes
\begin{equation}
\boxed{
\widehat\omega^2=\widehat q^4+\widehat q^2+\frac{\chi}{4}.
}
\label{eq:dimensionless_dispersion}
\end{equation}
so the same $\chi$ that controls static monotonic versus oscillatory response fixes the low-$q$ spectral gap. Figure~\ref{fig:dispersion} displays the corresponding dimensionless spectra.

\begin{figure}[t]
\includegraphics[width=0.94\linewidth]{dispersion.png}
\caption{Undamped dimensionless dispersion relation, Eq.~\eqref{eq:dimensionless_dispersion}, for the same three values of $\chi$ used in Fig.~\ref{fig:green}. The static-response classification and the dynamic spectral gap depend on the same constitutive parameters.}
\label{fig:dispersion}
\end{figure}

This provides the key falsification advantage of the construction. A finite-worldtube solver should not be allowed to fit a static ``memory'' profile with one set of $B,T,\kb$ and then use unrelated parameters for its oscillation spectrum. The same coefficients must account for both.

\section{Interpretation within the SAT$\rightarrow$\HsH{} sandbox}

The model suggests a precise replacement for a vague phrase such as ``the past exerts a force.'' A localized encounter changes the present extended configuration of a physical carrier; the present deformation then stores information about that encounter. In this language,
\begin{equation}
\boxed{
\text{interaction history}\;\longrightarrow\;\text{present geometric stress state}
}
\end{equation}
rather than
\begin{equation*}
\text{interaction history}\;\longrightarrow\;\text{new nonlocal force law}.
\end{equation*}

This distinction is important. An extended carrier can possess mechanical memory without requiring the equations of motion to contain an explicit time-history functional. Such a mechanism is familiar in ordinary extended media: a present configuration can encode previous loading. The speculative \HsH{} step is not the mathematics of elastic memory; it is the assertion that particle worldtubes possess analogous physical internal degrees of freedom.

The present calculation also separates three concepts that have sometimes been conflated in exploratory SAT language:
\begin{enumerate}[label=(\roman*)]
\item \textit{topological relation}: whether two carriers are linked, braided, or otherwise constrained;
\item \textit{local winding or deformation density}: how concentrated the relation is along the carrier;
\item \textit{constitutive persistence}: how $B$, $T$, and $\kb$ transmit or relax that deformation.
\end{enumerate}
Topological persistence does not guarantee a fixed force magnitude, and a mechanical backtug need not itself be topological.

In the current project taxonomy, ``Interbraid'' denotes direct carrier--carrier mechanics while ``Electrogravity'' denotes carrier--surface/medium response. The present model belongs to the first category. A later coupled model could let the shared surface or medium modify the restoring coefficient,
\begin{equation}
\kb\rightarrow \kb^{\rm IB}+\kb^{\rm EG},
\end{equation}
which would allow one interaction sector to move the other across the response boundary $\chi=1$. That possibility is a new conjecture, not a result of the present derivation.

Likewise, nothing here identifies the mechanical relation with quantum entanglement. A future claim of that kind would have to recover the relevant quantum-state structure and Bell-type correlations, not merely a persistent classical coupling.

\section{Solver protocol and tight predictions}

The model is deliberately easy to attack numerically. A useful finite-worldtube benchmark should implement the following protocol.

\subsection{Static test}

Construct two finite carriers with independently specified $B$, $T$, and $\kb$. Apply a localized relative displacement or force near $s=0$, remove the external drive, and relax the configuration. Sweep
\begin{equation}
10^{-2}\lesssim\chi\lesssim10^{2}.
\end{equation}
Without fitting the response class itself, test the following predictions:
\begin{align}
\chi<1 &: \quad \text{no sign-changing tail; sum of exponentials},\\
\chi=1 &: \quad \text{double-root critical profile},\\
\chi>1 &: \quad \text{damped sign-changing oscillatory tail}.
\end{align}
In the supercritical regime, fit $\alpha$ and $\beta$ from the profile and compare them directly with Eq.~\eqref{eq:alphabeta}.

\subsection{Dynamic test}

Using the same carrier and the same $B$, $T$, and $\kb$, excite small-amplitude modes over a range of $q$. Measure $\omega(q)$ and compare with Eq.~\eqref{eq:dispersion}. The stringent test is joint consistency: the coefficients inferred from the static Green function must predict the dynamic spectrum without re-fitting the constitutive law.

\subsection{Failure criteria}

The sandbox mechanism should be rejected or enlarged if any of the following occur in the regime where the linear assumptions are valid:
\begin{enumerate}[label=(\alph*)]
\item no persistent deformation remains after a localized encounter;
\item the static profile changes class at a parameter combination inconsistent with $T^2=4B\kb$;
\item the measured decay/oscillation constants disagree with Eqs.~\eqref{eq:yroots} or \eqref{eq:alphabeta};
\item the coefficients required by the static profile fail to predict the dynamic dispersion relation;
\item finite-core, nonlinear, topological, or reconnection effects dominate before the linear benchmark regime is reached.
\end{enumerate}
The last outcome is not a rescue clause. It identifies the specific additional physics that a more complete model must represent explicitly.

\section{Limitations and causal status}

Several boundaries should remain explicit.

First, a four-dimensional block representation does not itself imply mechanical coupling between separated temporal locations. Equation~\eqref{eq:energy} adds an extended-carrier constitutive assumption that standard point-particle mechanics does not contain.

Second, the coordinate $s$ is deliberately left as an intrinsic carrier coordinate. Interpreting it literally as timelike separation while simultaneously introducing an independent evolution parameter $\tau$ would require a careful causal construction. The present paper does not provide one. Accordingly, the static Green function should be read as a configuration-space relation, not evidence for controllable retrocausal signaling.

Third, the model is linear, local in $s$, and isotropic in the reduced deformation variable. Real finite worldtubes may possess self-contact, chirality, anisotropy, nonlinear tension, reconnection, multiple internal modes, and scale-dependent constitutive coefficients. Any of these can alter the response.

Fourth, the calculation does not derive a $1/r^2$ force law and does not reproduce general relativity. Therefore it is not presently a gravity model. Its value is more limited and more concrete: it supplies a mechanical null hypothesis for how an extended carrier could retain interaction history, together with a sharp constitutive discriminator.

\section{Conclusion}

If SAT is correct only in the modest sense that the standard four-dimensional map is a useful starting representation, then \HsH{} must earn every additional physical ingredient. A physical worldtube does not acquire history-dependent forces merely because its full trajectory can be drawn. The present sandbox model asks what happens if the tube is instead granted ordinary extended-medium degrees of freedom.

The minimal quadratic energy
\begin{equation*}
E=\frac12\int ds\,[B(u'')^2+T(u')^2+\kb u^2]
\end{equation*}
produces a complete and testable response structure. The dimensionless number
\begin{equation*}
\boxed{\chi=4B\kb/T^2}
\end{equation*}
separates monotonic stored deformation from damped spatial oscillation, with a critical boundary at $\chi=1$ or equivalently $\lI=2\lel$. The same coefficients predict the dynamic spectrum
\begin{equation*}
\boxed{\omega^2=(Bq^4+Tq^2+\kb)/\mu.}
\end{equation*}
Thus one candidate meaning of Interbraid ``backtug'' is simply the restoring stress associated with a deformation that remains encoded in the present extended carrier. In the supercritical regime, nominal backtug and forwardtug become alternating phases of the same damped spatial mode.

The result is not a new force. It is a challenge to the theory: if \HsH{} worldtubes are physical enough to store interaction history, their static and dynamic behavior should be mutually consistent and calculable. If they are not, the mechanism should be discarded.

\appendix

\section{Green-function derivation}

For $\chi<1$, the Fourier denominator factors as
\begin{equation}
Bq^4+Tq^2+\kb
=B(q^2+r_+^2)(q^2+r_-^2),
\end{equation}
where $r_\pm^2$ are given by Eq.~\eqref{eq:yroots}. Partial fractions give
\begin{equation}
\widetilde G(q)=\frac{1}{B(r_+^2-r_-^2)}
\left[
\frac{1}{q^2+r_-^2}-\frac{1}{q^2+r_+^2}
\right].
\end{equation}
Using
\begin{equation}
\int_{-\infty}^{\infty}\frac{dq}{2\pi}\frac{e^{iqs}}{q^2+a^2}
=\frac{e^{-a|s|}}{2a}
\end{equation}
produces Eq.~\eqref{eq:greensubcritical}.

At criticality the Fourier denominator becomes $B(q^2+r_0^2)^2$, and the standard inverse transform yields Eq.~\eqref{eq:greencritical}.

For $\chi>1$, the decaying characteristic roots are $\alpha\pm i\beta$. Evenness and continuity require the real-space form
\begin{equation}
G(s)=C e^{-\alpha|s|}
\left[
\cos(\beta|s|)+\frac{\alpha}{\beta}\sin(\beta|s|)
\right].
\end{equation}
Integrating Eq.~\eqref{eq:static} across $s=0$ gives the jump condition
\begin{equation}
B[G''']_{0^-}^{0^+}-T[G']_{0^-}^{0^+}=1.
\end{equation}
For the even solution $G'$ is continuous at the origin, while $[G''']=2G'''(0^+)$. This fixes
\begin{equation}
C=\frac{1}{4B\alpha(\alpha^2+\beta^2)},
\end{equation}
which gives Eq.~\eqref{eq:greensupercritical}.

\section{Reproducible dimensionless benchmark}

The figures use
\begin{equation}
B=T=\mu=1,
\qquad
\kb=\chi/4,
\end{equation}
with $\chi\in\{0.25,1,4\}$. The static profiles are numerical evaluations of
\begin{equation}
G(s)=\frac{1}{\pi}\int_0^\infty dq\,
\frac{\cos(qs)}{q^4+q^2+\chi/4}
\end{equation}
normalized by $G(0)$. The dynamic curves evaluate
\begin{equation}
\widehat\omega(q)=\sqrt{q^4+q^2+\chi/4}.
\end{equation}
No parameter was fitted to historical SAT constants or particle data.

\begin{thebibliography}{9}

\bibitem{gravitytwist}
N. McKnight, \emph{GRAVITY TWIST.txt}, SAT Theory Archive, archival research conversation/note, file path \path{2026/GRAVITY TWIST.txt}, Satobloc/SAT\_THEORY\_ARCHIVE\_2023-25 (accessed October 2, 2026), \url{https://github.com/Satobloc/SAT_THEORY_ARCHIVE_2023-25}.

\bibitem{forcesacrosstime}
N. McKnight, \emph{Forces Across Temporal Points - to 8-13-24}, archived \HsH{} conversation, file path \path{DEVELOPMENT_FULL_CONVOS/30SEP26_DUMP/002 Forces Across Temporal Points - to 8-13-24.txt}, Satobloc/HsH (accessed October 2, 2026), \url{https://github.com/Satobloc/HsH}.

\bibitem{landau}
L. D. Landau and E. M. Lifshitz, \emph{Theory of Elasticity}, 3rd ed., Course of Theoretical Physics Vol. 7 (Pergamon Press, Oxford, 1986).

\bibitem{audoly}
B. Audoly and Y. Pomeau, \emph{Elasticity and Geometry: From Hair Curls to the Non-linear Response of Shells} (Oxford University Press, Oxford, 2010).

\end{thebibliography}

\end{document}
